\section{实验：问题与验证对应}
在已有实机测试支持三个设计有效的基础上，实验按表示收益、反馈收益和任务表现组织，并用于核对理论条件与收益来源：首先在独立人工审核的分层视频子集上，对比自动标注迭代前后及质量诊断消融，报告人工时间、各类mask与几何误差，并单列接触、初现和重遮挡片段，检验式\eqref{eq:noise}及质量选择偏差；其次在相同示范、视角、种子和训练预算下，对比RGB ACT、语义ACT、无几何监督的同结构query及完整QToken，加入重复RGB的token预算对照，报告动作误差、几何读出和事件附近行为，以受控遮挡及query干预辅助检验式\eqref{eq:geobound}中的相关性与利用路径，事件真值独立审核，不能以同一预测mask自证任务完成；再次固定同一基础checkpoint、目标槽位和调度，比较零残差、仅计划上下文残差、最新观察MLP、无额外历史的同结构GRU及完整历史GRU，并用历史置乱、移除慢query和残差开关区分新观察、历史及语义贡献，报告式\eqref{eq:alignment}的实测两项、逐关节及计划年龄误差，不将不可直接观测的Bayes风险项当作已测量数值；随后注入可控延迟与丢帧，在相同设备负载下比较慢重规划、频繁完整ACT与快残差的时延分位数、目标槽覆盖率和恢复表现，检验式\eqref{eq:coverage}的适用范围；最后在两个真实平台的多类任务、配对初始条件和受控扰动下评价全系统成功率、事件到达率、拨空气次数、恢复率、完成后无效动作及任务耗时，统一超时规则并计入失败，使用多个训练种子和重复实机试验报告分层置信区间。跨A/B/C的全链路对比不能独立归因于快分支输入，相关消融须固定基础策略；未见对象/布局的划分应贯穿提示器、分割器和策略训练，模拟实验仅作为可选的隐状态与延迟控制验证，不替代实机结果。
